\input{devoir.sty}
% !TEX root = /home/eb/Dropbox/CPGE/Physique/Cours/[6]Pondes/[5]Dispersion/Interro.tex
      \begin{document}
      \EnteteInter{21/03/2024}{1}
      \begin{quest}
      \item Définir la vitesse de phase $v_{\vphi}(\omega)$ à partir d'une relation de dispersion $\udl{k}(\omega)$. Signification physique ?
      \ansbox{}
      \item Définir la vitesse de groupe $v_g(\omega)$ à partir d'une relation de dispersion $\udl{k}(\omega)$. Signification physique ?
      \ansbox{}
      \item Définir l'indice complexe $\udl{n}(\omega)$ à partir d'une relation de dispersion $\udl{k}(\omega)$. Signification physique des parties réelle et imaginaire ?
      \ansbox{}
      \item Quelles relations de passage s'appliquent aux ondes électromagnétiques à l'interface entre deux milieux ? En déduire l'expression des coefficients de réflexion $\udl{r}_E$ et de transmission $\udl{t}_E$ pour le champ électrique en fonction des indices $\udl{n}_1$ et $\udl{n}_2$ des deux milieux.
      \ansbox{}
      \end{quest}
      \FIN
      \end{document}
      
